\subsection{交换 Banach 代数的特征标}

定理 1. --- 设 A 是一个 Banach 代数，并设 \(\chi\colon\mathbf{A}\to\mathbf{C}\) 是一个代数态射 (参见 I, p.~9)。则 \(\chi\) 连续，且范数至多为 1。若 A 是幺的且 \(\chi\) 是一个保单位元态射，则 \(\chi\) 的范数为 1。

我们来证明 \(|\chi(x)| \leqslant \|x\|\)（对一切 \(x \in \mathbf{A}\)）。如有必要，把 A 换成由 \(x\) 生成的 Banach 代数，我们就可以假定 A 是交换的；这时 \(\chi \in \mathbf{X}'(\mathbf{A})\)。对每个 \(x \in \mathbf{A}\)，有 \(\chi(x) \in \mathrm{Sp}_{\mathrm{A}}'(x)\) (I, p.~9, 第 7 小节)，因而 \(|\chi(x)| \leqslant \varrho(x) \leqslant \|x\|\) (I, p.~28, 推论 5)，由此得第一个断言。

此外，若 A 是幺的，则等式 \(\chi(1)=1\) 蕴涵 \(\|\chi\|\geqslant1\)，由此得所要的等式。

注. --- 设 A 是一个交换 Banach 代数。由这个定理得出，\(X'(A)\) 上的逐点收敛拓扑与 \(X'(A)\) 上的诱导拓扑重合，后者由弱拓扑 \(\sigma(A', A)\)（A 的对偶 \(A'\) 的弱拓扑）诱导而来。

推论. --- 设 A 是一个交换 Banach 代数。空间 X'(A) 是紧的。空间 X(A) 是局部紧的，且若 A 有单位元，则它是紧的。

设 A' 是 Banach 空间 A 的对偶。它的单位球 \(A_{1}^{\prime}\) 是弱紧的 (EVT, III, p.~17, 推论 2)。由定理 1，我们有 \(\mathsf{X}^{\prime}(\mathsf{A}) \subset \mathsf{A}_{1}^{\prime}\)。此外，\(\mathsf{X}^{\prime}(\mathsf{A})\) 在 A' 中对于弱拓扑是闭的，因为它是下列弱闭集的交

\[
\mathrm{X} _ {x, y} = \left\{f \in \mathrm{A} ^ {\prime} \mid f (x y) - f (x) f (y) = 0 \right\} \quad (\text { for } x, y \in \mathrm{A}).
\]

由此得出 \(X'(A)\) 是紧的，且 \(X(A) = X'(A) - \{0\}\) 是局部紧的。

若 A 有单位元 1，则 \(\mathsf{X}(\mathsf{A})\) 是由 \(\chi \in \mathsf{X}'(\mathsf{A})\)（满足 \(\chi(1) = 1\)）所成的集合，因此是 \(\mathsf{X}'(\mathsf{A})\) 的一个闭子集，从而是紧的。

每个紧空间都同胚于某个适当的交换幺 Banach 代数 A 的 X(A) (参见 I, p.~32, 推论 2)。

定理 2. --- 设 A 是一个交换 Banach 代数。映射 \(\chi \mapsto \operatorname{Ker}(\chi)\) 是从 \(\mathsf{X}(\mathsf{A})\) 到 A 的正则极大理想所成的集合 \(\mathrm{J}(\mathrm{A})\) 上的一个双射。

设 I 是 A 的一个正则极大理想。它是闭的 (I, p.~22, 推论 2)，因而 A/I 是一个 Banach 代数。由于它是一个域 (A, VIII, p.~426, 命题 2)，Gelfand--Mazur 定理 (I, p.~26, 推论 2) 蕴涵 A/I 在 C 上维数为 1。于是 I 在 A 中的余维数为 1。定理于是由 I, p.~10 的引理 3 得出。

可能发生这样的情形：一个非零的、无单位元的交换 Banach 代数 A 没有极大理想；这时 \(X'(A)\) 缩为 \(\{0\}\) (参见 I, p.~186, 习题 31)。

这个定理表明，I, p.~11 中的集合 \(J(A)\) 与 \(\hat{A}\)（由于 A 是交换的，这两者视为同一）也可以与集合 \(X(A)\) 视为同一。因此有理由在 \(X(A)\) 上考虑弱拓扑和 Jacobson 拓扑，后者更粗 (I, p.~15, 命题 5)。Jacobson 拓扑的闭集 \(V(M)\)（对 \(M \subset A\)）是下列集合

\[
\{\chi \in \mathrm{X} (\mathrm{A}) \mid \chi (\mathrm{M}) = 0 \}
\]

（\(\mathsf{X}(\mathsf{A})\) 的子集），它们有时仍记作 \(\mathsf{V}(\mathsf{M})\)。类似地，对任一子集 \(\mathsf{M}\)（\(\mathsf{X}(\mathsf{A})\) 的子集），我们用 \(\Upsilon(\mathsf{M})\) 表示诸 \(\chi \in \mathsf{M}\) 的核之交这个理想；它等于 I, p.~13 中所定义的集合 \(\Upsilon(\mathsf{M}')\)，即对子集 \(\mathsf{M}' \subset \mathsf{J}(\mathsf{A})\)（当把 \(\mathsf{J}(\mathsf{A})\) 与 \(\mathsf{X}(\mathsf{A})\) 视为同一时与 M 对应的子集）所定义的那个集合。

当在 X(A) 中使用一个拓扑概念而没有指明所用的是哪个拓扑时，所说的始终是弱拓扑。

